% F082｜源结论：仅描述本条知识，不引入其他类型
\textbf{结论名称：}F082 双根号函数值域：几何意义法（距离差型）

\textbf{解决的问题：}两个平方根相减，求差的最大值、最小值、完整值域，并判断边界能否真正取得。

\textbf{条件：}设 $m,n\in\mathbb R$，$a,b>0$，$x\in\mathbb R$，
\[
 f(x)=\sqrt{(x-m)^2+a^2}-\sqrt{(x-n)^2+b^2}.
\]
取 $P(x,0)$、$A(m,a)$、$B(n,b)$，则 $f(x)=PA-PB$。

\textbf{第一步：}两段距离之差的绝对值不超过两定点之间的距离：
\[
 |f(x)|=|PA-PB|\leq AB
 =\sqrt{(m-n)^2+(a-b)^2}=:D.
\]
\textbf{注意：}这只是一个范围限制，不能据此认定正、负两个端点都能取到。

\textbf{完整值域：}令 $d=|m-n|$，则
\[
 \boxed{
 \begin{array}{c|c}
 \text{条件}& f(x)\text{ 的值域}\\\hline
 a>b&(-d,D]\\
 a<b&[-D,d)\\
 a=b,\ m\ne n&(-d,d)\\
 a=b,\ m=n&\{0\}
 \end{array}}
\]
当 $a\ne b$ 时，$A,B$ 连线延长后与 $x$ 轴交于 $P_0$，
\[
 \boxed{x_0=\frac{bm-an}{b-a}}.
\]
若 $a>b$，在 $x=x_0$ 处达到最大值 $D$；若 $a<b$，在 $x=x_0$ 处达到最小值 $-D$。

\textbf{与距离和型的联系：}原来上下不同侧的点，可以把其中一个关于 $x$ 轴翻折；每个根号不变，但求差时需要把两个固定点放在同侧，检查延长线能否碰到 $x$ 轴。若 $a=b$，两点连线平行于 $x$ 轴（或两点重合），就不能套交点公式。

\textbf{速答提醒：}求得边界后必须问“取等的点 $P$ 是否真的存在且在题目允许范围内”。若题目限制 $x$ 的范围，要重新检查。
